% SOURCE_AUTHORITY: sga5_fr_workpass.tex, lines 12835--12846 (LF-normalized unit).
% SOURCE_SHA256: 791F4EFFC5E02832D5D77ED03518C8156D6F07E4C8238B03545DB93D883FBB28
Comparando com (5.2.1), encontra-se facilmente $\tau=\lambda_n^*(\sigma_2)$. De fato, se $g\ne e$, a igualdade entre $\tau(g)$ e $\lambda_n^*(\sigma_2)$ é imediata por (4.5.2). Se $g=e$, a fórmula a verificar escreve-se
\[
\EP(C')-\operatorname{card}(Y')
=
N\EP(C)-\sum_{y\in Y}\sigma'_y(e),
\qquad N=\operatorname{card}(G),
\]
isto é, por (4.5.2),
\[
\EP(C')-N\EP(C)=-\deg(\mathfrak D_{C'/C}),
\]
que não é senão a fórmula clássica de Hurwitz.
\end{prooftext}
